\documentclass[14pt,a4paper]{extarticle}
\usepackage[T2A]{fontenc}
\usepackage[utf8]{inputenc}
\usepackage[english,russian]{babel}
\usepackage{lmodern}
\usepackage{geometry}
\usepackage{array}
\usepackage{tabularx}
\usepackage{longtable}
\usepackage{ltablex}
\usepackage{booktabs}
\usepackage{hyperref}
\usepackage{indentfirst}
\usepackage{microtype}
\usepackage{graphicx}
\usepackage{setspace}
\usepackage{listings}
\usepackage{graphicx}
\usepackage{caption}
\usepackage{float}
\usepackage{tocloft}
\usepackage{titlesec}
\usepackage{newtxtext,newtxmath}

\titleformat{\section}        % форматируем \section
  {\normalfont\bfseries}      % стиль: обычный + жирный
  {\thesection.}              % номер секции с точкой
  {2mm}                       % расстояние между номером и текстом
  {}                           % перед заголовком можно добавить что-то
\titlespacing*{\section}{1.5cm}{*3}{*1}  % {отступ слева}{отступ сверху}{отступ снизу}

\titleformat{\subsection}        % форматируем \section
  {\normalfont\bfseries}      % стиль: обычный + жирный
  {\thesubsection.}              % номер секции с точкой
  {2mm}                       % расстояние между номером и текстом
  {}                           % перед заголовком можно добавить что-то
\titlespacing*{\subsection}{1.5cm}{*3}{*1}  % {отступ слева}{отступ сверху}{отступ снизу}

\captionsetup{justification=centering}
\renewcommand{\ttdefault}{cmtt}
\keepXColumns
\setlength{\parindent}{3em}
\newgeometry{left=2.5cm,right=2.5cm,top=3cm,bottom=2.5cm}

\renewcommand{\contentsname}{\hspace*{4em}{\smallСодержание}}
\setlength{\cftbeforesecskip}{0.5em}
\renewcommand{\cftsecleader}{\cftdotfill{\cftdotsep}}
\renewcommand{\cftsecaftersnum}{.\hspace{2mm}}

\title{Отчёт по "Управление ИТ-проектами"}
\author{Выполнил студент группы М-РИТ-25-1, Станиславчук Сергей}

\begin{document}

\setstretch{1.5}
\noindent
\begin{minipage}[t]{0.01\textwidth}
    \includegraphics{~/programmer/++Programmer/5_1/LSTU.jpg}
\end{minipage}%
\begin{minipage}[c]{1.1\linewidth}
    \centering
    {\small
    МИНИСТЕРСТВО НАУКИ И ВЫСШЕГО ОБРАЗОВАНИЯ РФ \\[1em]
    ФЕДЕРАЛЬНОЕ ГОСУДАРСТВЕННОЕ БЮДЖЕТНОЕ \\
    ОБРАЗОВАТЕЛЬНОЕ УЧРЕЖДЕНИЕ ВЫСШЕГО ОБРАЗОВАНИЯ \\
    «ЛИПЕЦКИЙ ГОСУДАРСТВЕННЫЙ ТЕХНИЧЕСКИЙ УНИВЕРСИТЕТ»
    }
\end{minipage}%

\vspace{3em}

\noindent
\hspace{1em}%
\begin{minipage}[t]{0.8\textwidth}
\textbf{Институт} \hspace{2em} компьютерных наук \\
\textbf{Кафедра} \hspace{2em} автоматизированных систем управления
\end{minipage}


\vspace{6em}

\begin{center}
\textbf{Практическая работа №10} \\
по дисциплине \\ <<Управление ИТ-проектами>>
\end{center}

\vspace{6em}

\noindent
\begin{tabular*}{\textwidth}{@{\extracolsep{\fill}} l c r }
    Студент М-РИТ-25-1 & \makebox[4cm]{\hrulefill} & Станиславчук С.М. \\[2em]
Руководитель & \\
К.т.н., доцент & \makebox[4cm]{\hrulefill} & Ведищев В.В. \\
\end{tabular*}

\vfill
\begin{center}
    Липецк, 2026
\end{center}

\thispagestyle{empty}

\newpage
\tableofcontents
\newpage

\section{Примеры систем командной разработки}

Для организации совместной разработки программных продуктов используются специализированные системы, обеспечивающие хранение кода, управление изменениями и координацию участников проекта. К таким системам относятся GitLab, GitHub, Bitbucket, а также системы управления задачами, такие как Jira.

В рамках данного проекта основным инструментом является система контроля версий Git, размещённая на платформе GitHub. Она обеспечивает централизованное хранение исходного кода, управление ветками и поддержку совместной работы разработчиков.

\section{Организация взаимодействия}

Взаимодействие внутри команды организуется через систему контроля версий и систему управления задачами. Каждая задача фиксируется в трекере, после чего назначается ответственному исполнителю.

Разработка ведётся в отдельных ветках, соответствующих функциональности или исправлению ошибок. Бэкенд- и фронтенд-разработчики работают параллельно, согласуя интерфейсы взаимодействия. После завершения работы изменения проходят процедуру проверки, что позволяет выявить ошибки до интеграции в основную ветку.

Специалист по информационной безопасности участвует в проверке изменений, связанных с обработкой данных и механизмами защиты. Тестировщик проверяет реализованный функционал и формирует отчёты об ошибках. Системный администратор обеспечивает работоспособность инфраструктуры и участвует в развёртывании системы.

Таким образом, взаимодействие строится на принципах распределённой разработки с централизованным контролем изменений.

\section{Контроль версий}

Контроль версий осуществляется с использованием Git. Основная ветка проекта содержит стабильную версию системы, в то время как разработка ведётся в отдельных ветках, создаваемых для реализации конкретных задач.

Используется модель ветвления, включающая ветки \texttt{main}, \texttt{develop}, \texttt{feature/*}, \texttt{bugfix/*} и \texttt{hotfix/*}. Это позволяет изолировать изменения, упростить их тестирование и контролировать процесс интеграции.

Каждое изменение фиксируется в виде коммита с описанием выполненной работы. Перед объединением изменений в основную ветку проводится проверка кода, а также тестирование. Для обозначения стабильных версий используются теги, соответствующие принятой схеме версионирования.

Данный подход обеспечивает отслеживаемость изменений, возможность возврата к предыдущим версиям и повышение качества разработки программного продукта.
\end{document}
